\section{架构、调度与性能机制的统一关系}

\subsection{从网格到执行通道的完整映射}

从软件层次到硬件执行资源的完整映射链为

\[
\boxed{
\begin{array}{c}
\text{网格}\rightarrow\text{线程块}\rightarrow\text{流式多处理器驻留}\\[2mm]
\rightarrow\text{线程束}\rightarrow\text{处理块调度}\rightarrow\text{执行通道}
\end{array}
}
\]

线程块决定资源分配和协作范围；线程束决定调度与控制发散的粒度；处理块和执行资源决定线程束指令如何获得实际执行机会。把这三个粒度区分开，就能避免把“线程块同时驻留”“线程束被调度”和“线程在核心上执行”混为同一件事。

\subsection{四个性能概念之间的因果链}

\begin{enumerate}
  \item \textbf{资源限制决定驻留规模。} 线程数、线程块数、寄存器和共享内存等限制共同决定一个流式多处理器能容纳多少线程块与线程束。
  \item \textbf{驻留规模决定占用率。} 驻留线程束数与架构最大线程束数之比就是占用率。
  \item \textbf{占用率影响延迟隐藏机会。} 驻留线程束越多，某个线程束停顿时找到另一个就绪线程束的机会通常越高。
  \item \textbf{控制发散降低单个线程束的有效通道利用。} 即使有足够多驻留线程束，如果每个线程束内部大量通道因为分支而失活，执行资源的有效利用仍会下降。
\end{enumerate}

因此，线程块大小不是一个只影响“线程有多少”的参数。它同时改变每块线程束数、流式多处理器可驻留线程块数、资源装箱效率和可用于延迟隐藏的线程束池；如果算法中的控制结构又与线程索引相关，还要进一步考虑线程束内部的发散模式。

\subsection{概念边界总结}

\begin{center}
\small
\begin{tabularx}{\textwidth}{>{\bfseries}p{0.18\textwidth} p{0.22\textwidth} X}
\toprule
概念 & 主要粒度 & 回答的问题 \\
\midrule
线程块驻留 & 线程块/流式多处理器 & 一个线程块能否进入某个流式多处理器，以及同时能驻留多少块 \\
线程束调度 & 线程束/处理块 & 哪个已经驻留且就绪的线程束现在获得指令发射机会 \\
控制发散 & 单个线程束内部 & 同一线程束的活动线程是否沿不同控制路径执行 \\
延迟隐藏 & 多个驻留线程束 & 当前线程束等待时，是否有其他就绪线程束可以使用执行资源 \\
占用率 & 流式多处理器 & 驻留线程束池相对架构最大容量有多大 \\
线程块内同步 & 单个线程块 & 同一线程块中的线程如何在共享资源上建立协作和屏障 \\
跨线程块协作 & 多线程块/更高作用域 & 如何在不依赖普通线程块调度顺序的前提下建立受约束的协作 \\
\bottomrule
\end{tabularx}
\end{center}
